%\section{PAGINA 5}

\section{Análisis Usando el Teorema de Norton}

A continuación, se presentan los circuitos que utilizaremos:

\begin{figure}[htbp]
    \centering
    \begin{circuitikz}[european]
        \draw (0,0) node[below] {B} to[battery, l={$V_s=\SI{100}{\volt}$}, invert] (0,4) node[above] {A} to[R, l={$R_1=\SI{1}{\kilo\ohm}$}] (4,4) node[above] {D};
        \draw (4,4) to[battery, l={$V_2=\SI{150}{\volt}$}] (4,2) to[R, l={$R_2=\SI{2.2}{\kilo\ohm}$}] (4,0) node[below] {C};
        \draw (4,4) -- (8,4) node[above] {F};
        \draw (8,4) to[R, l={$R_3=\SI{1.2}{\kilo\ohm}$}] (8,2) to[ammeter, l=$I_2$] (8,0) node[below] {G};
        \draw (8,0) -- (4,0) -- (0,0); % Conecta G -> C -> B
        \fill (0,0) circle (1.5pt); % Nodo B
        \fill (0,4) circle (1.5pt); % Nodo A
        \fill (4,4) circle (1.5pt); % Nodo D
        \fill (4,0) circle (1.5pt); % Nodo C
        \fill (8,4) circle (1.5pt); % Nodo F
        \fill (8,0) circle (1.5pt); % Nodo G
    \end{circuitikz}
    \caption{Solución de una Red Compleja por el Teorema de Norton.}
    \label{fig:circuito original}
\end{figure}

\begin{figure}[htbp]
    \centering
    \begin{circuitikz}[european]
        \draw (0,0) node[below] {B} to[battery, l={$V_s=\SI{100}{\volt}$}, invert] (0,4) node[above] {A} to[R, l={$R_1=\SI{1}{\kilo\ohm}$}] (4,4) node[above] {D};
        \draw (4,4) to[battery, l={$V_2=\SI{150}{\volt}$}] (4,2) to[R, l={$R_2=\SI{2.2}{\kilo\ohm}$}] (4,0) node[below] {C};
        \draw (4,4) -- (8,4) node[above] {F};
        \draw (8,4) to[ammeter, l=$I_N$] (8,0) node[below] {G}; % R3 y I2 reemplazados por I_N
        \draw (8,0) -- (4,0) -- (0,0); % Conecta G -> C -> B
        \fill (0,0) circle (1.5pt); % Nodo B
        \fill (0,4) circle (1.5pt); % Nodo A
        \fill (4,4) circle (1.5pt); % Nodo D
        \fill (4,0) circle (1.5pt); % Nodo C
        \fill (8,4) circle (1.5pt); % Nodo F
        \fill (8,0) circle (1.5pt); % Nodo G
    \end{circuitikz}
    \caption{Cálculo de la corriente de Norton ($I_N$) cortocircuitando los terminales F y G.}
    \label{fig:circuito-in}
\end{figure}

\begin{figure}[htbp]
    \centering
    \begin{circuitikz}[european]
        \draw (0,0) node[below] {B} to[short] (0,4) node[above] {A} to[R, l={$R_1=\SI{1}{\kilo\ohm}$}] (4,4) node[above] {D};
        \draw (4,4) to[R, l={$R_2=\SI{2.2}{\kilo\ohm}$}] (4,0) node[below] {C};
        \draw (4,4) -- (8,4) node[above] {F};
        \draw (8,0) -- (4,0) -- (0,0); % Conecta G -> C -> B
        \fill (0,0) circle (1.5pt); % Nodo B
        \fill (0,4) circle (1.5pt); % Nodo A
        \fill (4,4) circle (1.5pt); % Nodo D
        \fill (4,0) circle (1.5pt); % Nodo C
        \fill (8,4) circle (1.5pt); % Nodo F
        \node at (8,0) [below] {G}; % Nodo G como terminal abierto
    \end{circuitikz}
    \caption{Cálculo de la resistencia de Norton ($R_N$) vista desde los terminales F y G.}
    \label{fig:circuito-rn}
\end{figure}

\begin{figure}[htbp]
    \centering
    \begin{circuitikz}[european]
        \draw (0,0) node[below] {B} to[isource, l={$I_N=\SI{31.8}{\milli\ampere}$}] (0,4) node[above] {A} to[short] (4,4) node[above] {D};
        \draw (4,4) to[R, l={$R_N=\SI{687.5}{\ohm}$}] (4,0) node[below] {C};
        \draw (4,4) -- (8,4) node[above] {F};
        \draw (8,4) to[R, l={$R_3=\SI{1.2}{\kilo\ohm}$}] (8,2) to[ammeter, l=$I_N$] (8,0) node[below] {G};
        \draw (8,0) -- (4,0) -- (0,0); % Conecta G -> C -> B
        \fill (0,0) circle (1.5pt); % Nodo B
        \fill (0,4) circle (1.5pt); % Nodo A
        \fill (4,4) circle (1.5pt); % Nodo D
        \fill (4,0) circle (1.5pt); % Nodo C
        \fill (8,4) circle (1.5pt); % Nodo F
        \fill (8,0) circle (1.5pt); % Nodo G
    \end{circuitikz}
    \caption{Circuito Equivalente de Norton con la carga conectada.}
    \label{fig:circuito-equivalente-norton}
\end{figure}

\noindent En el lazo 1,

\begin{equation}
\begin{split}
    &100 - I_4 \times (1000) = 0 \\
    &I_4 = \frac{100}{1000} = \SI{0.1}{\ampere}
\end{split} \tag{27-16}
\end{equation}

En el bucle 2, 

\begin{equation}
\begin{split}
    &2200 \times I_5 + 150 = 0 \\
    &I_5 = \frac{150}{2200} = \SI{0.0682}{\ampere}
\end{split} \tag{27-17}
\end{equation}

La figura 27-5b muestra que $I_N=I_4-I_5$, puesto que los sentidos de 
$I_4$ y $I_5$ son opuestos. Por lo tanto,

\begin{equation}
\begin{split}
    &I_N = I_4 - I_5 = \SI{0.1}{\ampere} - \SI{0.0682}{\ampere} = \SI{0.0318}{\ampere}
\end{split} \tag{27-18}
\end{equation}  

La resistencia del generador de Norton $R_N$ es la resistencia total de 
la red original, es decir, la resistencia total de la red original es 
$R_N=\SI{687.5}{\ohm}$. Por lo tanto, el generador de Norton 
es $R_N=\SI{687.5}{\ohm}$ y $I_N=\SI{0.0318}{\ampere}$.

Nuevamente hallamos que el valor de la corriente en $R_3$ es el mismo que 
el valor obtenido utilizando las leyes de Kirchhoff y el teorema de 
Thévenin.

\section{Resumen}
\begin{itemize}[series=resumenlist]
\item El teorema de Norton ofrece otro método de resolver circuitos complicados. 
Por el teorema de Norton se puede reemplazar una red complicada de dos 
terminales por un circuito equivalente sencillo que actúa como el circuito 
original en la carga conectada a los dos terminales.
\item El circuito equivalente se compone de un generador de corriente, llamado
generador Norton ($I_N$) en paralelo con la resistencia interna $R_N$ del 
generador de Norton y en paralelo con los cuales está también conectada 
la carga $R_L$. La corriente $I_N$ del generador se divide entonces entre 
$R_N$ y $R_L$. La figura 27-1b muestra el generador de Norton y su carga. 
Es el equivalente de Norton de un circuito complicado tal como el representado 
en forma de bloques en la figura 27-1a.
\item Para determinar la corriente de Norton $I_N$, se cortocircuita la carga 
y se calcula la corriente que fluye por el circuito. Esta corriente de cortocircuito 
es $I_N$. Para calcular $I_N$ puede ser necesario hacer uso de las leyes de Ohm 
y de Kirchhoff.
\item Para determinar la resistencia de Norton $R_N$ que actúa en paraleo con 
el generador de Norton se aplica la misma técnica que para hallar la 
la resistencia de Thévenin en el experimento precedente. Se procedde como 
sigue: Abrir la carga en los dos terminales en cuestión de la red original. Reemplazar 
todas la fuentes de tensión por sus resistencias internas. Luego calcular $R_N$, 
resistencia en los terminales abiertos de la carga, hacia el circuito.
\item El teorema de Norton se aplica a un circuito con una o mas fuentes de 
alimentación.
\item Una vez reemplazada la red complicada por el circuito equivalente de Norton 
se puede hallar la corriente $I_L$ que pasa por la carga aplicando la fórmula 
siguiente (fig. 27-1b):

\begin{equation}
\begin{split}
    &I_L = \frac{I_N \times R_N}{R_N + R_L}
\end{split} \tag{}
\end{equation}
\end{itemize}